\begin{lemma}[Perfect sub-super-sequences give perfect sub-multi-sequences]
Let $h$ be locally constant and let $h^{\dcl}$ be its induced
super-sequence on $F^h$.  If a restriction of $h^{\dcl}$ to a sub-front is
perfect, then $h$ has a perfect sub-multi-sequence.
\end{lemma}
\begin{proof}
The structure statement \noderef{Front} gives
$Y.\mathsf{Infinite}$ from the field $hF'.\mathsf{infinite\_base}$.
Apply \noderef{InfiniteSetEnumeration} to choose an increasing enumeration
$z\in\mathsf{IncSeq}$ with

\[
\operatorname{range}(z)=Y.
\tag{1}
\]

We prove that the fixed-left restriction is perfect.  Fix an arbitrary
$x\in\mathsf{IncSeq}$ and put

\[
a=z\circ x,\qquad b=z\circ\mathsf{ShiftMap}(x).
\tag{2}
\]

Here both $a$ and $b$ are in $\mathsf{IncSeq}$ by the composition statement
\noderef{IncSeqComp} and the definition of $\mathsf{ShiftMap}$ in
\noderef{ShiftMap}.  Thus their ranges are infinite by the enumeration
meaning of \noderef{IncSeq}.  Moreover,
$\operatorname{range}(a)\subseteq Y$: if $n=a(k)$, then the composition
equation from \noderef{IncSeqComp} gives $n=z(x(k))$, so (1) gives $n\in Y$.
The identical argument gives
$\operatorname{range}(b)\subseteq Y$, since $b(k)=z(\mathsf{ShiftMap}(x)(k))$.

The density field of the Front statement \noderef{Front}, applied first to
$\operatorname{range}(a)\subseteq Y$ and its infinitude, supplies
$s\in F'$ with

\[
s\in F'\quad\text{and}\quad
\mathsf{ProperPrefixSet}(s,\operatorname{range}(a)).
\tag{3}
\]

Applying the same density field to $\operatorname{range}(b)$ supplies
$t\in F'$ with

\[
t\in F'\quad\text{and}\quad
\mathsf{ProperPrefixSet}(t,\operatorname{range}(b)).
\tag{4}
\]

The explicit hypothesis $hsub$ is the subfamily inclusion
$F'\subseteq\mathsf{DecidingFrontSet}(h)$, so (3) and (4) also give the
deciding-front memberships of the corresponding subtypes.

The coherence equation for the value map on the deciding front is the
property supplied by \noderef{DecidingValue}; for the particular map $f$
in this theorem it is explicitly the hypothesis $hf$.  Therefore (3), with
the subtype $\langle s,hsub(s\in F')\rangle$, gives

\[
f.value\langle s,hsub(s\in F')\rangle=h(a),
\tag{5}
\]

and (4) gives, with the analogous subtype,

\[
f.value\langle t,hsub(t\in F')\rangle=h(b).
\tag{6}
\]

We next verify the finite-shift premise needed for perfectness.  The
restriction-shift identity \noderef{RestrictionShift}, instantiated at
$z$ and $x$, says

\[
b=z\circ\mathsf{ShiftMap}(x)
 =\mathsf{ShiftMap}(z\circ x)=\mathsf{ShiftMap}(a).
\tag{7}
\]

Consequently (4) is also
$\mathsf{ProperPrefixSet}(t,\operatorname{range}(\mathsf{ShiftMap}(a)))$.
Taking $a$ itself as the increasing-sequence witness, (3) and this last
equality are exactly the two conjuncts in the definition
\noderef{FiniteShift}; hence

\[
\mathsf{FiniteShift}(s,t).
\tag{8}
\]

Now instantiate the perfectness statement \noderef{PerfectSuperSequence}
for the restricted super-sequence
$\mathsf{SuperSequenceRestrict}(f,F',Y,hF',hsub)$, with the members $s,t$,
their membership proofs from (3) and (4), and the witness (8).  This gives

\[
r\bigl(\mathsf{SuperSequenceRestrict}(f,F',Y,hF',hsub).value
       \langle s,s\in F'\rangle,
       \mathsf{SuperSequenceRestrict}(f,F',Y,hF',hsub).value
       \langle t,t\in F'\rangle\bigr).
\tag{9}
\]

By the definition \noderef{SuperSequenceRestrict}, the first value in (9)
is $f.value\langle s,hsub(s\in F')\rangle$ and the second is
$f.value\langle t,hsub(t\in F')\rangle$.  Rewriting with (5) and (6)
turns (9) into

\[
r(h(a),h(b)).
\tag{10}
\]

Finally, \noderef{MultiSequenceRestrict} says that
$(\mathsf{MultiSequenceRestrict}(h,z))(x)=h(a)$ and, using (2),
$(\mathsf{MultiSequenceRestrict}(h,z))(\mathsf{ShiftMap}(x))=h(b)$.
Thus (10) is precisely the required relation for this arbitrary $x$.
Since $x$ was arbitrary, the defining condition in
\noderef{PerfectMultiSequence} proves that
$\mathsf{MultiSequenceRestrict}(h,z)$ is perfect.
\end{proof}
